\documentclass[12pt]{article}
\usepackage[utf8]{inputenc}
\usepackage[T1]{fontenc}
\usepackage{hyperref}
\title{Unggumi Language Resources for K-12 Curriculum Development}
\author{Compiled by Claude Sonnet 5.5 (claude-sonnet-5-5) for Andr\'as Kornai}
\date{}
\begin{document}
\maketitle

\section{Introduction}
Unggumi (ISO 639-3 xgu; Glottocode ungg1243) is a Western Worrorran language of the Kimberley region of northern Western Australia. Sources disagree on its status: Glottolog lists it as extinct with a comment that it is dormant, and ASJP gives extinct since 1995. It was split from Worrorra in ISO 639-3. I found no school curriculum or teaching materials.

\section{Primary Educational Resources}
I found no primers, textbooks or school programmes for Unggumi.

\section{Linguistic and Reference Materials}
The Glottolog entry is at \url{https://glottolog.org/resource/languoid/id/ungg1243}. The ASJP word-list entry is at \url{https://asjp.clld.org/languages/UNGGUMI}. The Wikipedia overview of the Worrorran family is at \url{https://en.wikipedia.org/wiki/Worrorran_languages}. Capell's 1940 comparative vocabulary includes Unggumi, but I did not see a direct link to it.

\section{Religious and Community Resources}
I found no scripture or community organisation resources specific to Unggumi.

\section{Digital and Audio Resources}
I found no audio, video or software resources.

\section{Curriculum Development Context}
Documentation appears limited to word lists and comparative work. A curriculum would have to draw on archival material and on closely related Worrorran languages, such as Worrorra, which is documented more fully. That should be done only with the descendant community's agreement and with Unggumi material kept separate from the related languages.
\end{document}